\section{Reciprocal depth budgets for the three actual arms}
\label{rdc-section}

Retain the actual primitive unramified root and quotient domain of
Lemma~\ref{sac-quotients}: $Q\ge7$, $3\nmid Q$, $n\ge5$ and
$\gcd(n,6)=1$. The three nonzero integer quotients $D_i$ are
pairwise coprime. Use its notation $T_n,T_1,c,t,c_1,\Delta$ and,
for $Y\ge5$, put
\[
 S_i=\sum_{q>Y}v_q(D_i)\log q,\qquad
 R_i=\sum_{\substack{q>Y\\q\mid D_i}}\log q,\qquad
 L_i=\sum_{q\le Y}v_q(D_i)\log q,\qquad
 L=\sum_{q\le Y}v_q(T_n)\log q.
\]
All sums are over primes. The actual height relations give
\begin{equation}\label{rdc-height-interval}
 t-\Delta-\log c_1-L_i\le S_i\le t,\qquad
 L_i\ge0,\quad \sum_iL_i\le L.
\end{equation}

Choose caps $h_i\in\mathbb Z_{>0}\cup\{\infty\}$. An infinite cap
places no depth restriction on that quotient. Let $I$ consist of the
finite-cap indices and define
\begin{align*}
 \beta_i&=3/h_i\quad(i\in I),& \beta_i&=0\quad(i\notin I),\\
 E_i^{(h_i)}&=\sum_{q>Y}(v_q(D_i)-h_i)_+\log q\quad(i\in I),&&\\
 C_h&=\sum_i\beta_i,& A_h&=\sum_i(\beta_i-1)_+,\qquad
 B_h=\max_i(\beta_i-1)_+.
\end{align*}
Thus $0\le A_h\le6$ and $0\le B_h\le2$, even for caps that vary
with the root and exponent.

\begin{theorem}[Exact finite reciprocal-cap compensation]\label{rdc-finite}
On this actual domain,
\begin{align}
 \log W_Y(T_n)\le{}&(3-C_h)t+A_h(\Delta+\log c_1)
                    +B_hL+\log T_1\notag\\
                 &+\sum_{i\in I}\beta_iE_i^{(h_i)}
                    -3\sum_{i\notin I}R_i.
                    \label{rdc-finite-bound}
\end{align}
The complete depths of every infinite-cap arm are included and its
negative radical contribution remains on the right.
\end{theorem}
\begin{proof}
At a supported prime, $e\le h_i+(e-h_i)_+$, so
$3R_i\ge\beta_i(S_i-E_i^{(h_i)})$ for $i\in I$.
Pairwise coprimality gives exact signed additivity on the quotient
product. Hence
\[
 \log W_Y(D_1D_2D_3)
 \le\sum_i(1-\beta_i)S_i+
       \sum_{i\in I}\beta_iE_i^{(h_i)}-3\sum_{i\notin I}R_i.
\]
When $\beta_i\le1$, use the upper bound $S_i\le t$. When
$\beta_i>1$, its coefficient is negative and
\eqref{rdc-height-interval} bounds its contribution by
\[
 (1-\beta_i)t+(\beta_i-1)(\Delta+\log c_1+L_i).
\]
Summation yields $(3-C_h)t+A_h(\Delta+\log c_1)+B_hL$.
Finally use $T_n=T_1|D_1D_2D_3|$ and the general overlap inequality
$\log W_Y(UV)\le\log W_Y(V)+\log|U|$ proved in
Theorem~\ref{sac-finite}. No coprimality of $T_1$ and the quotient
product is assumed.
\end{proof}

Write
\[
 \mathcal E_h=\sum_{i\in I}\beta_iE_i^{(h_i)}-
                       3\sum_{i\notin I}R_i.
\]
\begin{theorem}[Uniform signed-tail subclasses]\label{rdc-uniform}
With the absolute effective constants of Theorem~\ref{sac-subclass},
\begin{equation}\label{rdc-normalized}
 \frac{\log W_Y(T_n)}t
 \le3-C_h+\epsilon_h(n,Y)+\frac{\mathcal E_h}{t},
\end{equation}
where
\[
 \epsilon_h(n,Y)=A_hA_1\frac{\log(4n)}n+\frac{2A_h+3}{n}
       +B_hA_2\frac{Y^3\log Y\log^2(4n)}n.
\]
For $Y=n^{1/6}$, this error tends to zero uniformly in all actual
moving roots and caps. In particular, if
\begin{equation}\label{rdc-capacity}
 \sum_{i\in I}\frac1{h_i}\ge1
\end{equation}
and the actual finite-cap arms satisfy $v_q(D_i)\le h_i$ at every
$q>Y$, then $(\log W_Y(T_n))_+/t\to0$ uniformly as $n\to\infty$.
The same conclusion follows under the more general condition
$(\mathcal E_h)_+/t\to0$. If $C_h\ge3+\delta$ for a fixed
$\delta>0$ and the caps are met, then
\[
 \limsup\frac{\log W_Y(T_n)}t\le-\delta.
\]
\end{theorem}
\begin{proof}
Sector normalization gives the same unit-residual representation
with $\lambda=1/n$ as in Theorem~\ref{sac-subclass}. Thus its
all-prime two-place bounds control $\Delta/t$ and $L/t$ as stated
there. The remaining numerator obeys
\[
 A_h\log c_1+\log T_1
 \le\left(\frac{A_h}{2}+\frac32\right)\log Q
                          +A_h\log(2/\sqrt3).
\]
Divide by $t\ge n\log Q/2$, use $Q\ge7$ and
$2\log(2/\sqrt3)<\log7$, and obtain the bound $(2A_h+3)/n$.
Substitution proves \eqref{rdc-normalized}. The bounds $A_h\le6$,
$B_h\le2$ give uniform convergence of the error. Under the depth
caps all finite excesses vanish, so $\mathcal E_h\le0$.
Condition~\eqref{rdc-capacity} is exactly $C_h\ge3$.
The positive-part and fixed negative-margin conclusions follow.
\end{proof}

The earlier absorption therefore gives the restricted ABC estimate
for these specified subclasses, uniformly for sufficiently large $n$
for each fixed target epsilon. Bounded indices with arbitrarily
moving roots remain outside this argument. The signed estimate is
one-sided, rather than a claim that negative contributions disappear.

\begin{proposition}[The five maximal depth-cap patterns]\label{rdc-five}
Up to permutation, the maximal integer-cap patterns satisfying
\eqref{rdc-capacity}, under componentwise enlargement, are
\begin{equation}\label{rdc-patterns}
 (1,\infty,\infty),\quad(2,2,\infty),\quad
 (2,3,6),\quad(2,4,4),\quad(3,3,3).
\end{equation}
\end{proposition}
\begin{proof}
Sort $h_1\le h_2\le h_3$, with infinity largest. If $h_1=1$ the
triple is below the first pattern. If $h_1=h_2=2$, it is below the
second. If $h_1=2,h_2=3$, the reciprocal condition forces $h_3\le6$.
If $h_1=2,h_2\ge4$, it forces $h_2=h_3=4$.
If $h_1=3$, the other two must also be three; $h_1\ge4$ is
impossible. Each displayed pattern has reciprocal sum exactly one.
\end{proof}

The first pattern is a new useful form of the compensation mechanism:
one quotient that is squarefree above $Y$ pays for the full depths
in both other arms. More precisely, at that pattern
\begin{equation}\label{rdc-one-arm}
 \log W_Y(T_n)\le2\Delta+2\log c_1+2L+\log T_1
                 +3\bigl(E_i^{(1)}-R_j-R_k\bigr).
\end{equation}
Thus squarefreeness can be weakened to
$(E_i^{(1)}-R_j-R_k)_+/t\to0$. The pattern $(2,2,\infty)$
recovers the preceding SA criterion. Patterns $(2,3,6)$ and $(2,4,4)$
permit positive signed contributions at depth above three, while
another actual arm compensates them; both have $A_h=B_h=1/2$.
No automatic arithmetic membership is asserted for these conditions.

\begin{proposition}[Sharpness only for the abstract finite ledger]
\label{rdc-abstract-sharpness}
When fixed caps have $C_h<3$, their depth restrictions and equal
total arm masses alone do not force a sublinear signed integer-weight
ledger. This statement does not give an actual Eisenstein counterexample.
\end{proposition}
\begin{proof}
Let $H$ be the least common multiple of the finite caps, or one if
there are none, and set $t_m=mH$. For a finite cap $h_i$, take a
one-entry list of depth $h_i$ and integer weight $t_m/h_i$.
For an infinite cap take depth $t_m$ and weight one. All finite caps
hold and all three masses $S_i$ equal $t_m$. If $r$ caps are infinite,
the exact signed sum is
\[
 \sum_i(S_i-3R_i)=(3-C_h)t_m-3r.
\]
Its ratio to $t_m$ tends to $3-C_h>0$. These finite lists are not
asserted to be prime factorizations or to satisfy the actual additive
arm relation and common norm. They show that extra arithmetic input
is needed beyond this ledger; they exclude no actual norm route.
\end{proof}

\paragraph{Formal verification and unresolved scope.}
The separate \texttt{ReciprocalDepthArithmetic} module proves seventeen
declarations using finite natural depths and nonnegative integer
weights. It proves arbitrary-cap depth/radical inequalities, their
common-denominator interface $b_ih_i=3d$, the negative-coefficient
height costs, the canonical three-arm budget and its uniform bounds,
and exact one-, two- and three-finite-cap classifications after
clearing positive denominators. The fresh Lean~4.32.0 build checks all
seventeen new declarations and seventy-one dependency declarations
with warnings treated as errors and only the standard three axioms.
This module does not define actual prime factorizations or real
logarithms, formalize the analytic estimates, or prove cap membership.

The unresolved arithmetic object is the actual excess-minus-radical
term $\mathcal E_h$, constrained by the common power, additive arm
relation and norm. Large private depths can invalidate every simple
cap pattern without refuting its net-credit sufficient condition.
Neither the abstract sharpness family nor the proved sufficient
classes settle the general signed tail or ABC.
